
Multi-dimensional trustworthiness evaluation frameworks for large language models (LLMs) assess dimensions such as truthfulness, robustness, fairness, safety, and privacy independently, yet deployment failures frequently span multiple dimensions simultaneously. This study investigates whether trustworthiness dimension failures exhibit coupling patterns---systematic co-occurrence---that are invisible to independent evaluation. Through phi coefficient analysis on synthetic benchmark data emulating production evaluation structure (n=1500 instances across 5 dimensions and 3 simulated model profiles), we demonstrate methodology for detecting and characterizing coupling when present. Results from synthetic validation show coupling is detectable (phi 0.33-0.40, p < 1e-13) but sparse: no model exhibits coupling in three or more dimension pairs after Bonferroni correction. Two dominant patterns emerge---truthfulness-robustness coupling (phi 0.36-0.40) and fairness-safety coupling (phi 0.33-0.40)---across GPT-4, Claude-3, and Llama-3 profiles. Partial correlation analysis controlling for instance difficulty yields retention of 86-161\%, establishing that coupling in synthetic data reflects designed correlation structure rather than spurious artifacts of hard instances failing uniformly. Models exhibit observationally distinct coupling profiles in synthetic data (Mantel r < 0.7), though statistical confirmation requires larger samples than the proof-of-concept provides (n=100 per dimension; n $\geq$ 500 required). All findings derive from synthetic coupling data designed to validate measurement methodology; real-world coupling magnitudes in production LLMs remain uncharacterized pending access to gated benchmarks (MultiTrust, TrustLLM). If validated on real models, sparse coupling detection would enable coupling-aware model selection for deployments requiring multiple trustworthiness guarantees and reveal vulnerability clusters requiring targeted rather than universal mitigation strategies.
